% Gesamttest «Lineare Funktionen» — Quelle des PDFs (python3 scripts/build-lp-pdf.py).
% Musterlösung und Raster: bewertungspaket.tex. Alle Werte mit python3 nachgerechnet (03.10.2026).
\documentclass[11pt]{article}
\usepackage{../lp-druck}
\renewcommand{\lpname}{Lineare Funktionen}
\renewcommand{\lpteilgebiet}{3.2}
\renewcommand{\lpfuss}{Gesamttest · 22 Punkte}
\begin{document}
\lpkopf{Gesamttest}{%
  \vspace{4pt}\noindent\begin{tabularx}{\linewidth}{@{}X X@{}}
  Code (kein Name): \hrulefill & Punkte: \hrulefill\ / 22
  \end{tabularx}}

\begin{lpinfo}{Anleitung}
\textbf{Zeit:} rund 20 Minuten. \textbf{Hilfsmittel:} keine — alle drei Kompetenzen des Teilgebiets 3.2 tragen im Lehrplan den Vermerk «auch ohne Hilfsmittel».
Schreib zu jeder Aufgabe den \textbf{Rechenweg} auf — bewertet wird der Weg.
Danach: Lösung scannen oder fotografieren (ohne Namen und Standort) und mit dem \emph{Bewertungspaket} bewerten lassen.
\end{lpinfo}

\lpteil{Teil A · Gerade und Koeffizienten lesen}{ohne Taschenrechner · 9 Punkte}

\begin{lpaufgabe}{G1}{Graph zeichnen}{3 P}
\begin{minipage}[t]{0.52\linewidth}
Zeichne den Graphen von $f(x) = -1.5x + 3$ in das Koordinatensystem (Wertetabelle erlaubt) und gib die Nullstelle an.
\lpplatz{5}
\end{minipage}\hfill
\begin{minipage}[t]{0.42\linewidth}\vspace{0pt}\centering
\begin{tikzpicture}
\begin{axis}[width=6.4cm,height=6.4cm,axis lines=middle,xmin=-3,xmax=4,ymin=-3,ymax=4,
  xtick={-3,...,4},ytick={-3,...,4},grid=both,grid style={lplinie},tick label style={font=\scriptsize},
  xlabel=$x$,ylabel=$y$,axis equal image,
  yticklabel style={fill=white,inner sep=1pt,xshift=-2pt},xticklabel style={fill=white,inner sep=1pt}]
\end{axis}
\end{tikzpicture}
\end{minipage}
\end{lpaufgabe}

\begin{lpaufgabe}{G2}{Vom Graphen zur Gleichung}{3 P}
\begin{minipage}[t]{0.52\linewidth}
Gib die Gleichung der abgebildeten Geraden an. Die markierten Punkte liegen auf Gitterpunkten.
Berechne damit $f(-3)$.
\lpplatz{5}
\end{minipage}\hfill
\begin{minipage}[t]{0.42\linewidth}\vspace{0pt}\centering
\begin{tikzpicture}
\begin{axis}[width=6.4cm,height=6.4cm,axis lines=middle,xmin=-4,xmax=4,ymin=-5,ymax=3,
  xtick={-4,...,4},ytick={-5,...,3},grid=both,grid style={lplinie},tick label style={font=\scriptsize},
  xlabel=$x$,ylabel=$y$,axis equal image,
  xlabel style={anchor=north west},   % sonst liegt das «x» auf der Geraden und auf (3 | 0)
  yticklabel style={fill=white,inner sep=1pt,xshift=-2pt},xticklabel style={fill=white,inner sep=1pt}]
\addplot[lpblau,very thick,domain=-4:4] {2/3*x-2};
\addplot[only marks,mark=*,color=lpgruen] coordinates {(0,-2) (3,0) (-3,-4)};
\end{axis}
\end{tikzpicture}
\end{minipage}
\end{lpaufgabe}

\begin{lpaufgabe}{G3}{Einsetzen und deuten}{3 P}
Gegeben ist $f(x) = 0.5x - 3$.
\begin{enumerate}[label=(\alph*),leftmargin=*,itemsep=1pt]
\item Gib den $y$-Achsenabschnitt und die Nullstelle an.
\item Liegt $P(10 \mid 2)$ auf dem Graphen? Rechne nach.
\item Was bedeutet $m = 0.5$ für den Verlauf der Geraden? Ein Satz genügt.
\end{enumerate}
\lpplatz{7}
\end{lpaufgabe}

\lpteil{Teil B · Typen und Lage}{ohne Taschenrechner · 5 Punkte}

\begin{lpaufgabe}{G4}{Zu einer Geraden}{3 P}
Gegeben ist $g: y = -4x + 1$.
\begin{enumerate}[label=(\alph*),leftmargin=*,itemsep=1pt]
\item Stelle die Gleichung der Geraden auf, die parallel zu $g$ ist und durch $(0 \mid -3)$ geht.
\item Welche Steigung hat eine Gerade, die senkrecht auf $g$ steht?
\item Warum ist $x = 5$ keine Funktion?
\end{enumerate}
\lpplatz{6}
\end{lpaufgabe}

\begin{lpaufgabe}{G5}{Drei Fälle}{2 P}
Ordne zu: Welche der drei Gleichungen $y = 0.5x$, $y = -2$ und $x = 4$ beschreibt eine
\emph{proportionale} Funktion, welche eine \emph{konstante} Funktion, und welche \emph{keine} Funktion?
Begründe den letzten Fall in einem Satz.
\lpplatz{5}
\end{lpaufgabe}

\lpteil{Teil C · Gleichung aufstellen}{ohne Taschenrechner · 8 Punkte}

\begin{lpaufgabe}{G6}{Zwei Punkte}{3 P}
Eine Gerade geht durch $A(-3 \mid 6)$ und $B(2 \mid -4)$. Stelle ihre Funktionsgleichung auf.
Welcher besondere Typ liegt hier vor?
\lpplatz{6}
\end{lpaufgabe}

\begin{lpaufgabe}{G7}{Steigung und Punkt}{2 P}
Eine Gerade hat die Steigung $m = -0.5$ und geht durch $P(6 \mid 1)$. Stelle ihre Funktionsgleichung auf.
\lpplatz{5}
\end{lpaufgabe}

\begin{lpaufgabe}{G8}{Tank}{3 P}
Ein Tank enthält 300 Liter. Ab dem Zeitpunkt $t = 0$ fliessen gleichmässig 12 Liter pro Minute ab.
Stelle $V(t)$ auf ($t$ in Minuten, $V$ in Litern), berechne, wann der Tank leer ist, und gib die
Definitionsmenge im Sachzusammenhang an.
\lpplatz{7}
\end{lpaufgabe}
\end{document}
